\section{Introduction}
\label{sec:introducao}

With the expansion of digital platforms, there is an increase in the volume of data generated and processed as a result of the digitization of daily activities and the growing need to collect and process large datasets~\cite{sqlBigData:2022}. In this context, the JavaScript Object Notation (JSON) format has emerged as a popular choice for representing and exchanging data in web applications due to its flexibility and ease of use. However, with its popularity, significant challenges have arisen, such as data parsing~\cite{misonJson:2017}, performance issues~\cite{dataInterchangeMobile:2011}, and concerns about the quality and efficiency of the application~\cite{pageExp:2022}. The increase in data volume implies higher computational costs. For this reason, it is crucial that developers understand the impact of their technological choices and coding practices on computational resource usage and the energy consumption of applications~\cite{pwaEnergy:2022}.


Previous studies compare data interaction formats in applications, addressing performance, computational costs, and environmental concerns. For instance, \citet{dataInterchangeMobile:2011} compare processing and energy consumption costs in mobile applications using XML, JSON, and Protocol Buffers.  \citet{dataConsumptionIot:2021} show that the choice of data encoding technology has a significant impact on performance and resource efficiency. Additionally, \citet{jsonXmlCaseStudy:2009} compare the performance of JSON and XML in a client-server environment, measuring transmission speed and computing resource usage. However, other factors such as the different processing approaches of data formats, both on the client and server sides, and their impact on energy consumption, performance, and user experience in data exchange for web applications must be considered.


Processing large volumes of JSON data has become a bottleneck in terms of performance and cost~\cite{misonJson:2017}, which highlights the importance of understanding and analyzing the cost of data interaction. Efficiently addressing this issue not only seeks to reduce the costs of handling data~\cite{misonJson:2017}, but also contributes to alleviating environmental concerns by reducing energy consumption and Carbon Dioxide (\ce{$CO_2$}) emissions. According to \citet{mvcEnergy:2023}, in the year 2020, the Information and Communication Technology (ICT) sector was responsible for 1.43 billion tons of \ce{$CO_2$}, representing 2.5\% of global greenhouse gas emissions.


In this scenario, the goal of this paper is to compare the efficiency and performance between Client-Side Rendering (CSR) and Server-Side Rendering (SSR) in web applications. To achieve this goal, the following specific objectives were proposed: (i) analyze performance aspects during data transactions; (ii) identify and compare the computational costs involved in data exchanges; and (iii) investigate and measure application performance and interactivity based on Google's Core Web Vitals (CWV) metrics.

%JP: adicionei uma breve descrição da metodologia adotada, sugiro uma revisão para verificar se está coerente. Também seria interessante adicionar mais detalhes da base de dados utilizada, como vi que faltam informações referentes a esse ponto no texto achei melhor não adicionar por agora.
To accomplish these objectives, a systematic methodology was adopted. A back-end application, implemented using Node.js, was made responsible for handling two parameters: data size and rendering type. A front-end application, developed in JavaScript, was designed to request the data, render the content, and display the results. Performance evaluation was automated through a custom script. This script leveraged the Lighthouse library for comprehensive auditing, \texttt{chrome-launcher} for initiating clean Chrome browser instances, and Puppeteer for collecting profiling and computational cost data.

The results demonstrate statistically significant differences between CSR and SSR across various metrics. It was observed that SSR rendering performed better in metrics such as Time to Interactive (TTI) and Largest Contentful Paint (LCP), while the CSR approach excelled in First Contentful Paint (FCP), and Speed Index (SI). Based on this data, it is expected that software engineers will be able to determine the most advantageous scenarios for each approach, taking into account response time, memory consumption, and CWV indicators. Thus, this study provides a guide to assist in choosing the most appropriate rendering method, contributing to the development of more efficient and user-friendly web applications that optimize performance and minimize computational costs.


This study is structured as follows: Section~\ref{sec:background} presents the theoretical background related to data interaction in web applications. Section~\ref{sec:related_work} discusses related work. Section~\ref{sec:materiais} details the materials and methods used in this research. Sections~\ref{sec:results} and~\ref{sec:discussion} present and discuss the results. Threats to validity are addressed in Section~\ref{sec:threats_to_validity}. Finally, Section~\ref{sec:conclusion} presents the conclusions.
